%\needspace{4\baselineskip}
\color{uniblue}{\subsection{Защита от неполнофазного режима}}
\color{black}

\begin{enumerate}[label=\arabic{section}.\arabic{subsection}.\arabic{enumi}, labelsep=4pt, leftmargin=0pt, itemindent=57pt, itemsep=0pt, parsep=0pt] % если что можно крутить itemsep


\item
Защита предназначена для ликвидации неполнофазного режима работы выключателей с 
пофазным приводом, возникающего в результате возможного отказа переключения одной или двух фаз 
выключателя при коммутации. 

\item
Ввод в работу ЗНР осуществляется с помощью уставки \textbf{«Работа ЗНР»}. Оперативный 
ввод/вывод ЗНР  производится с помощью БУ (блока управления функциями) \textbf{«Вывод ЗНР »}. Пуск ЗНР 
формируется при условии появления конфигурируемых сигналов \textbf{«ЗНФ В1»} или \textbf{«ЗНФ В2»} и 
отключенном положении смежного выключателя (для схем с двумя выключателями), с контролем 
протекания тока нулевой или обратной последовательности. Выбор количества контролируемых 
выключателей задается уставкой \textbf{«ЗНР контр. В2»}, \textbf{«ЗНР контр. В3»}. Порог срабатывания ЗНР по току 
нулевой последовательности задается уставкой \textbf{«Iср ЗНР»}. Выдержка времени срабатывания ЗНР 
задается уставкой \textbf{«Тср ЗНР »}. 

\begin{figure}[H]
\centering
\includegraphics[width=1\textwidth,height=1\textheight,keepaspectratio]{\fbpath/271.0221 ЗНР/img/ZNR.pdf}
\captionof{figure}{Функциональный блок ЗНР}\label{pict:ZNR}
\end{figure}

\small
\setlength{\extrarowheight}{0.05cm} % высота ячеек в таблице
\begin{longtable}
{|>{\centering\arraybackslash}m{4.5cm}|>{\centering\arraybackslash}m{4.5cm}|>{\centering\arraybackslash}m{3cm}|>{\centering\arraybackslash}m{3.5cm}|}
\caption{Перечень уставочных параметров функции ЗНР\hfill\vspace{-0.5\baselineskip}}\label{ksn:tbl2} \\ 
\hhline{|-|-|-|-|}
\rowcolor{gray!30}
Наименование уставки & Диапазон & Шаг &  Значение по умолчанию \\ 
\hhline{|-|-|-|-|}
\endfirsthead

\caption*{\hspace{3pt}\emph{Продолжение таблицы \ref{ksn:tbl2}\hfill\vspace{-0.5\baselineskip}}} \\ % сделано по ГОСТ 2.105 п.6.8.7
\hhline{|-|-|-|-|}
\rowcolor{gray!30}
Наименование уставки & Диапазон & Шаг &  Значение по умолчанию \\  
\hhline{|-|-|-|-|}
\endhead
\hhline{|-|-|-|-|}

\endfoot
\endlastfoot

%%%%%%%%%%%%содержимое%%%%%%%%%%%%%%

\end{longtable}
\setlength{\extrarowheight}{0.0cm} % возврат параметра высота ячеек в таблице
\normalsize


\end{enumerate} 